% !TEX root = ../thuthesis-main.tex

\chapter{总结与展望}

\section{工作总结}

本文围绕数字电路时序图的自动化解析问题，提出了一种基于多模态大语言模型的两阶段解析框架，实现了从时序图图像到WaveJSON结构化表示、再到可综合Verilog代码的完整处理流水线。

在数据构建方面，本文针对该领域缺乏标注数据的核心瓶颈，设计了一套以Verilog仿真为基础的自动化数据合成链路。通过引入大语言模型替代随机激励生成具有功能覆盖性的测试平台，使合成的时序图能够充分展示模块的关键工作状态；同时利用随机管线的产出作为低成本探针对候选模块进行质量预选，并采用基于并查集的防泄漏划分策略确保评估的可靠性，最终从108{,}971个原始Verilog源文件中筛选出5{,}847个高质量训练样本。

在评价体系方面，本文提出了面向WaveJSON结构特性的信号级F1评价指标，通过信号名匹配、保持符号展开与时钟等价性处理，解决了传统文本相似度指标对信号顺序敏感且无法区分语义错误严重程度的问题。该指标同时作为强化学习的奖励函数发挥了双重作用。

在模型训练方面，本文将视觉理解与代码推理解耦为两个独立阶段：第一阶段使用Qwen3-VL-4B-Instruct通过LoRA微调实现时序图到WaveJSON的转换，在117个测试样本上达到99.1\%的解析成功率与0.969的平均F1；第二阶段使用Qwen3-8B通过LoRA微调实现WaveJSON到Verilog代码的转换，以8B参数量超越了参数量高出数十倍的开源通用模型。在此基础上，本文引入GRPO强化学习以仿真对拍的F1分数作为连续值奖励信号，将仿真通过率从71.8\%提升至83.8\%，整体F1从0.648提升至0.692，与参数量远超一个数量级的Claude闭源模型（0.751）差距收窄至0.059，充分验证了面向特定任务的数据合成与强化学习训练流水线的有效性。


\section{局限性与未来展望}

本文方法仍存在若干局限性，同时也指明了未来的研究方向。

其一，训练与评测数据均为合成生成的WaveDrom风格时序图，模型对真实工程场景中的时序图（如芯片数据手册中的协议规范图、EDA工具导出的仿真波形截图或手绘草图）的泛化能力尚未验证。未来可从数据手册、IP规格文档等来源采集真实时序图数据，结合多种渲染风格与分辨率进行跨域训练，弥合合成数据与实际应用之间的分布差异。

其二，两阶段流水线存在误差传播的固有风险——第一阶段的信号名错误或波形长度偏差会直接劣化第二阶段的代码质量。未来可探索端到端的单阶段架构，或引入跨阶段反馈机制（如利用第二阶段的编译错误信息引导第一阶段修正输出）以缓解误差累积。

其三，基于强化学习的优化仍处于初步探索阶段，且训练尚不充分。当前的探索主要验证了仿真奖励作为优化信号的可行性，但模型的改善集中在代码结构规范性上，功能逻辑层面的深层能力仍有较大提升空间。未来可在奖励函数设计（如结合形式化验证工具提供更强的正确性保证）、训练策略与规模等方面进行更深入的探索。

其四，当前数据合成依赖单一的Verilog源代码库，模块类型以中小规模RTL设计为主。未来可将数据来源扩展到工业级IP设计，覆盖更复杂的参数化结构与总线协议交互，并将本文方法应用于设计文档与RTL实现的一致性检查等实际工程场景，最终集成到EDA工具链中构建智能设计辅助系统。
